\chapter{Chiral Hochschild Cohomology and Koszul Duality}

\section{Motivation: The Deformation Problem for Chiral Algebras}

\subsection{Historical Genesis and Physical Motivation}

The development of Hochschild cohomology for chiral algebras emerged from three independent streams of thought that converged in the 1990s. First, physicists studying marginal deformations of conformal field theories needed to understand when a perturbation $S \to S + \lambda \int \phi(z,\bar{z}) d^2z$ preserves conformal invariance. Seiberg \cite{Sei88} recognized that exactly marginal deformations correspond to closed elements in a certain cohomology theory. Second, mathematicians following Gerstenhaber's deformation theory \cite{Ger63} sought to extend Hochschild cohomology to vertex algebras. Third, Beilinson-Drinfeld's formalization of chiral algebras \cite{BD04} as factorization algebras demanded a cohomology theory respecting the geometric structure.

The fundamental question is: Given a chiral algebra $\mathcal{A}$ on a smooth curve $X$, what are its infinitesimal deformations that preserve the chiral structure? In classical algebra, if we deform an associative multiplication $\mu: A \otimes A \to A$ to $\mu_t = \mu + t\phi$, the associativity constraint
\[
\mu_t(\mu_t \otimes \text{id}) = \mu_t(\text{id} \otimes \mu_t)
\]
must hold to first order in $t$. Expanding, we find $\phi$ must satisfy
\[
\mu \circ (\phi \otimes \mathrm{id}) + \phi \circ (\mu \otimes \mathrm{id}) - \mu \circ (\mathrm{id} \otimes \phi) - \phi \circ (\mathrm{id} \otimes \mu) = 0
\]
which in terms of elements reads $a\phi(b,c) - \phi(ab,c) + \phi(a,bc) - \phi(a,b)c = 0$.
This is precisely the Hochschild 2-cocycle condition $\delta^2\phi = 0$. The obstruction to extending to second order lives in $HH^3(A,A)$.

For chiral algebras, the situation is far richer. A deformation must preserve:
\begin{enumerate}
\item The $\mathcal{D}_X$-module structure encoding locality
\item The chiral multiplication $\mu: j_*j^*(\mathcal{A} \boxtimes \mathcal{A}) \to \Delta_*\mathcal{A}$
\item The singularity structure along the diagonal
\item The operator product expansion coefficients
\end{enumerate}

\subsection{Why Configuration Spaces Enter}

The appearance of configuration spaces is not a mathematical convenience but a physical necessity. In quantum field theory, the principle of locality states that operators commute at spacelike separation. On a curve $X$, this means the commutator $[\phi_1(z_1), \phi_2(z_2)]$ must vanish for $z_1 \neq z_2$. All nontrivial structure is thus encoded in the approach $z_1 \to z_2$.

The configuration space $C_n(X) = \{(z_1,\ldots,z_n) \in X^n : z_i \neq z_j\}$ parametrizes positions where operators don't collide. Its compactification $\overline{C}_n(X)$ adds boundary divisors $D_{ij} = \{z_i = z_j\}$ that encode collision limits. A deformation of the chiral algebra must specify how the algebraic structure changes as points approach these divisors.

\section{Construction of the Chiral Hochschild Complex}

\subsection{The Cochain Spaces}

\begin{definition}[Chiral Hochschild Complex - Geometric Realization]
For a chiral algebra $\mathcal{A}$ on a smooth curve $X$, define the degree $n$ cochains as
\[
C^n_{\text{chiral}}(\mathcal{A}) = \Gamma\left(\overline{C}_{n+2}(X), j_*j^*\mathcal{A}^{\boxtimes (n+2)} \otimes \Omega^n_{\overline{C}_{n+2}(X)}(\log D)\right)
\]
where:
\begin{itemize}
\item $\overline{C}_{n+2}(X)$ is the Fulton-MacPherson compactification
\item $j: C_{n+2}(X) \to \overline{C}_{n+2}(X)$ is the open embedding
\item $\mathcal{A}^{\boxtimes (n+2)}$ denotes the external tensor product on $X^{n+2}$
\item $\Omega^n_{\overline{C}_{n+2}(X)}(\log D)$ are $n$-forms with logarithmic poles along the boundary divisor $D$
\end{itemize}
\end{definition}

The index $n+2$ (rather than $n$) appears because Hochschild cohomology involves one output, $n$ inputs, and one evaluation point. Explicitly, a degree $n$ cochain is a sum of expressions
\[
\phi = \sum_I a_0^{(I)}(z_0) \otimes a_1^{(I)}(z_1) \otimes \cdots \otimes a_n^{(I)}(z_n) \otimes a_{\infty}^{(I)}(z_{\infty}) \otimes \omega_I
\]
where $a_i^{(I)} \in \mathcal{A}$ and $\omega_I$ is an $n$-form on $\overline{C}_{n+2}(X)$ with logarithmic singularities.

\subsection{The Differential: Three Components United}

The differential $d: C^n_{\text{chiral}} \to C^{n+1}_{\text{chiral}}$ has three components reflecting the algebraic, geometric, and operadic structures:

\begin{theorem}[The Chiral Hochschild Differential; \ClaimStatusProvedHere]
The differential decomposes as
\[
d = d_{\text{int}} + d_{\text{fact}} + d_{\text{config}}
\]
where:
\begin{enumerate}
\item $d_{\text{int}}$: internal differential from the $\mathcal{D}_X$-module structure
\item $d_{\text{fact}}$: factorization using chiral multiplication  
\item $d_{\text{config}}$: de Rham differential on configuration space
\end{enumerate}
\end{theorem}

\begin{proof}
We verify $d^2 = 0$ by analyzing all nine combinations:

\textbf{Pure terms:}
\begin{align}
d_{\text{int}}^2 &= 0 \quad \text{($\mathcal{A}$ is a complex of $\mathcal{D}_X$-modules)} \\
d_{\text{config}}^2 &= 0 \quad \text{(de Rham differential squares to zero)} \\
d_{\text{fact}}^2 &= 0 \quad \text{(associativity of chiral multiplication)}
\end{align}

\textbf{Mixed terms:} The crucial cancellation 
\[
d_{\text{fact}} \circ d_{\text{config}} + d_{\text{config}} \circ d_{\text{fact}} = 0
\]
follows from the Arnold-Orlik-Solomon relations. For any configuration of three points:
\[
d\log(z_1-z_2) \wedge d\log(z_2-z_3) + d\log(z_2-z_3) \wedge d\log(z_3-z_1) + d\log(z_3-z_1) \wedge d\log(z_1-z_2) = 0
\]

This relation, discovered by Arnold \cite{Arn69} in studying configuration spaces of hyperplanes and generalized by Orlik-Solomon \cite{OS80}, encodes the fact that three points on a curve have only two degrees of freedom. Geometrically, it says the sum of exterior derivatives around a triangle vanishes.

The remaining mixed terms vanish because $d_{\text{int}}$ commutes with both other differentials by $\mathcal{D}_X$-linearity.
\end{proof}

\subsection{Explicit Formula for the Differential}

For a cochain $\phi \in C^n_{\text{chiral}}$, the differential acts by:

\begin{align}
(d_{\text{int}}\phi)(&z_0,\ldots,z_{n+1}) = \sum_{i=0}^{n+1} (-1)^i d_{\mathcal{A}}(\phi(z_0,\ldots,\hat{z}_i,\ldots,z_{n+1})) \\
(d_{\text{fact}}\phi)(&z_0,\ldots,z_{n+1}) = \sum_{i=1}^n (-1)^i \text{Res}_{z_i = z_0} \phi(\mu(z_0,z_i),z_1,\ldots,\hat{z}_i,\ldots,z_{n+1}) \\
&+ \sum_{1 \leq i < j \leq n} (-1)^{i+j} \phi(z_0,\ldots,\mu(z_i,z_j),\ldots,\hat{z}_i,\ldots,\hat{z}_j,\ldots,z_{n+1}) \\
(d_{\text{config}}\phi)(&z_0,\ldots,z_{n+1}) = d_{\overline{C}_{n+2}}(\phi)
\end{align}

where $\hat{z}_i$ denotes omission and $\mu$ is the chiral multiplication.

\section{Computing Cohomology via Bar-Cobar Resolution}

\subsection{The Resolution Strategy}

Computing Hochschild cohomology directly from the definition is typically intractable. The bar-cobar resolution provides a systematic approach:

\begin{theorem}[Hochschild via Bar-Cobar; \ClaimStatusProvedHere]
For any chiral algebra $\mathcal{A}$, there is a quasi-isomorphism
\[
C^{\bullet}_{\text{chiral}}(\mathcal{A}) \simeq \text{Hom}_{\text{ChirAlg}}(\Omega^{\text{ch}}(\overline{B}^{\text{ch}}(\mathcal{A})), \mathcal{A})
\]
where $\Omega^{\text{ch}}(\overline{B}^{\text{ch}}(\mathcal{A}))$ is the cobar construction of the bar complex.
\end{theorem}

\begin{proof}
The proof has three steps:

\textbf{Step 1: Bar gives cofree resolution.} 
The geometric bar complex $\overline{B}^{\text{ch}}(\mathcal{A})$ constructed in Chapter 4 is a cofree chiral coalgebra resolving $\mathcal{A}$:
\[
\overline{B}^{\text{ch}}(\mathcal{A}) \xrightarrow{\epsilon} \mathcal{A}
\]

\textbf{Step 2: Cobar gives free resolution.}
Applying the cobar functor (Chapter 5) yields a free chiral algebra resolution:
\[
\Omega^{\text{ch}}(\overline{B}^{\text{ch}}(\mathcal{A})) \xrightarrow{\eta} \mathcal{A}
\]

\textbf{Step 3: Hom computes Ext.}
By definition, 
\[
\text{Ext}^n_{\text{ChirAlg}}(\mathcal{A}, \mathcal{A}) = H^n(\text{Hom}_{\text{ChirAlg}}(\Omega^{\text{ch}}(\overline{B}^{\text{ch}}(\mathcal{A})), \mathcal{A}))
\]
The left side is precisely $HH^n_{\text{chiral}}(\mathcal{A})$ by definition.
\end{proof}

\subsection{The Spectral Sequence}

The double complex structure induces a spectral sequence:

\begin{theorem}[Hochschild Spectral Sequence; \ClaimStatusProvedHere]
There exists a spectral sequence
\[
E_1^{n,q} = H^q(\overline{C}_{n+2}(X), \mathcal{A}^{\boxtimes (n+2)} \otimes \Omega^n_{\log}) \Rightarrow HH^{n+q}_{\text{chiral}}(\mathcal{A})
\]
where $n$ is the bar degree (number of internal insertions) and $q$ is the sheaf cohomology degree, which are independent indices.
\end{theorem}

For formal chiral algebras (quasi-isomorphic to their cohomology), this spectral sequence degenerates at $E_1$, giving:
\[
HH^k_{\text{chiral}}(\mathcal{A}) \cong \bigoplus_{n+q=k} H^q(\overline{C}_{n+2}(X), \mathcal{A}^{\boxtimes (n+2)} \otimes \Omega^n_{\log})
\]

\section{Koszul Duality for Chiral Algebras}

\subsection{Quadratic Chiral Algebras and Their Duals}

\begin{definition}[Quadratic Chiral Algebra]
A chiral algebra $\mathcal{A}$ is \emph{quadratic} if it admits a presentation
\[
\mathcal{A} = T_{\text{chiral}}(\mathcal{V})/(R)
\]
where:
\begin{itemize}
\item $\mathcal{V}$ is a locally free $\mathcal{O}_X$-module of generators
\item $T_{\text{chiral}}(\mathcal{V})$ is the free chiral algebra on $\mathcal{V}$
\item $R \subset j_*j^*(\mathcal{V} \boxtimes \mathcal{V})$ consists of quadratic relations
\end{itemize}
\end{definition}

The free chiral algebra requires care to define. Following Beilinson-Drinfeld:

\begin{definition}[Free Chiral Algebra]
The free chiral algebra on $\mathcal{V}$ is
\[
T_{\text{chiral}}(\mathcal{V}) = \bigoplus_{n \geq 0} \pi_{n*}\left(j_*j^*\mathcal{V}^{\boxtimes n} \otimes \mathcal{D}_{C_n(X)/X}\right)^{\Sigma_n}
\]
where $\pi_n: C_n(X) \to X$ is the projection and $\mathcal{D}_{C_n(X)/X}$ denotes relative differential operators.
\end{definition}

\begin{definition}[Koszul Dual]
The Koszul dual of a quadratic chiral algebra $\mathcal{A}$ is
\[
\mathcal{A}^! = T_{\text{chiral}}(\mathcal{V}^*)/(R^{\perp})
\]
where:
\begin{itemize}
\item $\mathcal{V}^* = \mathcal{H}om_{\mathcal{O}_X}(\mathcal{V}, \omega_X)$ is the dual shifted by the canonical bundle
\item $R^{\perp}$ consists of relations orthogonal to $R$ under the canonical pairing
\[
\langle \cdot, \cdot \rangle: j_*j^*(\mathcal{V}^* \boxtimes \mathcal{V}) \to j_*\omega_{X^2 \setminus \Delta}
\]
\end{itemize}
\end{definition}

\begin{remark}[What This Definition Actually Says]\label{rem:koszul-dual-meaning}
The Koszul dual $\mathcal{A}^!$ defined above is \emph{precisely the coalgebra} that bar constructs from $\mathcal{A}$. More precisely:

\begin{enumerate}
\item \textbf{Generator duality}: The generators of $\mathcal{A}^!$ are the duals of the generators of $\mathcal{A}$:
   $$\mathcal{V}^* = \mathcal{H}om_{\mathcal{O}_X}(\mathcal{V}, \omega_X)$$
   This means: if $\mathcal{A}$ has generators $\phi_1, \ldots, \phi_n$, then $\mathcal{A}^!$ has dual generators $\phi_1^*, \ldots, \phi_n^*$

\item \textbf{Relation orthogonality}: The relations $R^{\perp}$ in $\mathcal{A}^!$ are orthogonal to the relations $R$ in $\mathcal{A}$ under the residue pairing:
   $$\langle r, r^* \rangle = \int_{X^2 \setminus \Delta} r \wedge r^* = 0 \quad \text{for all } r \in R, r^* \in R^{\perp}$$
   This means: what is a relation in $\mathcal{A}$ becomes "freedom" in $\mathcal{A}^!$, and vice versa

\item \textbf{Bar computes this dual}: The bar construction $\bar{B}^{\text{ch}}(\mathcal{A})$ naturally produces a coalgebra whose generators are $\mathcal{V}^*$ and whose coproduct encodes the relations $R^{\perp}$
\end{enumerate}

Therefore, saying $\bar{B}^{\text{ch}}(\mathcal{A}) \simeq \mathcal{A}^!$ is not a new condition but rather a \emph{verification} that the bar construction does what we expect: it produces the Koszul dual coalgebra.
\end{remark}

\begin{example}[Explicit Correspondence for Heisenberg]
For the Heisenberg chiral algebra $\mathcal{H}$ with generator $\alpha(z)$ and OPE:
$$\alpha(z)\alpha(w) \sim \frac{1}{(z-w)^2}$$

The Koszul dual $\mathcal{H}^!$ has:
\begin{itemize}
\item \textbf{Dual generator}: $\alpha^*(z)$ with $\langle \alpha, \alpha^* \rangle = 1$ under residue pairing
\item \textbf{Coproduct}: 
   $$\Delta(\alpha^*) = \alpha^* \otimes 1 + 1 \otimes \alpha^* + \text{(higher order terms)}$$
   encoding the dual of the commutative algebra structure
\item \textbf{Bar construction}: $\bar{B}^{\text{ch}}(\mathcal{H})$ consists of forms like:
   $$\alpha(z_1) \otimes \cdots \otimes \alpha(z_n) \otimes \eta_{12} \wedge \eta_{23} \wedge \cdots$$
   whose residues extract the coproduct coefficients
\end{itemize}

The cobar $\Omega^{\text{ch}}(\mathcal{H}^!)$ reconstructs a commutative chiral algebra from this coalgebraic data.
\end{example}

\begin{remark}[Why "Koszul Dual" vs "Dual"?]
The term "Koszul dual" (rather than just "dual") emphasizes that:
\begin{enumerate}
\item This is a derived/homotopical notion (quasi-isomorphisms, not isomorphisms)
\item It involves a specific homological construction (bar-cobar)
\item It generalizes the classical Koszul duality for quadratic algebras
\item The duality is self-inverse: $(\mathcal{A}^!)^! \simeq \mathcal{A}$
\end{enumerate}

When $\mathcal{A}$ is quadratic, $\mathcal{A}^!$ recovers the classical quadratic dual. For non-quadratic chiral algebras, $\mathcal{A}^!$ is defined by the bar construction but maintains all the essential dualities.
\end{remark}

\subsection{The Universal Twisting Morphism}

The relationship between a chiral algebra and its Koszul dual is mediated by:

\begin{definition}[Universal Twisting Morphism]
A \emph{twisting morphism} $\tau: \mathcal{A}^! \to \mathcal{A}$ is a degree 1 map satisfying the Maurer-Cartan equation
\[
\partial \tau + \tau \star \tau = 0
\]
where $\star$ denotes convolution in $\text{Hom}(\overline{B}^{\text{ch}}(\mathcal{A}^!), \Omega^{\text{ch}}(\overline{B}^{\text{ch}}(\mathcal{A})))$.
\end{definition}

\begin{theorem}[Existence and Uniqueness; \ClaimStatusProvedElsewhere]
For a Koszul pair $(\mathcal{A}, \mathcal{A}^!)$, there exists a unique universal twisting morphism $\tau: \mathcal{A}^! \to \mathcal{A}$ that induces quasi-isomorphisms:
\begin{align}
\mathcal{A}^!_{\tau} &\simeq \overline{B}^{\text{ch}}(\mathcal{A}) \\
\mathcal{A}_{\tau} &\simeq \Omega^{\text{ch}}(\mathcal{A}^!)
\end{align}
where the subscript denotes twisting by $\tau$.
\end{theorem}

\begin{remark}[What the Twisting Morphism Does]\label{rem:twisting-morphism-meaning}
The twisting morphism $\tau$ is the \textbf{explicit map implementing the bar-cobar isomorphism}. Concretely:

\begin{enumerate}
\item \textbf{Direction}: $\tau: \mathcal{A}^! \to \mathcal{A}$ is a map from the dual coalgebra to the original algebra

\item \textbf{Maurer-Cartan equation}: The condition $\partial \tau + \tau \star \tau = 0$ ensures that $\tau$ intertwines the coalgebra differential on $\mathcal{A}^!$ with the algebra differential on $\mathcal{A}$

\item \textbf{Twisted structures}: 
   \begin{itemize}
   \item $\mathcal{A}^!_{\tau}$ is $\mathcal{A}^!$ with differential twisted by $\tau$
   \item $\mathcal{A}_{\tau}$ is $\mathcal{A}$ with structure twisted by $\tau$
   \item The theorem says these twisted structures are quasi-isomorphic to bar and cobar
   \end{itemize}

\item \textbf{Universality}: $\tau$ is universal in that any other twisting factors through it
\end{enumerate}

Geometrically, $\tau$ is realized by:
$$\tau(c) = \int_{C_2(X)} \text{ev}^* c \wedge K_{\text{twist}}$$
where $K_{\text{twist}}$ is a universal integration kernel on the configuration space.
\end{remark}

\begin{example}[Twisting for Fermion-Boson Duality]
For the Koszul pair (free fermions $\mathcal{F}$, $\beta\gamma$ system $\mathcal{BG}$):

The twisting morphism $\tau: \mathcal{F}^! \to \mathcal{F}$ is given by:
$$\tau(\psi^*)(z) = \int_{\mathbb{C}} \psi(w) \cdot \frac{dw}{(z-w)}$$

This map:
\begin{itemize}
\item Takes the dual generator $\psi^*$ of $\mathcal{F}^!$
\item Integrates it against the fermion field $\psi$ with the basic kernel $\frac{1}{z-w}$
\item Produces a twisted field that satisfies bosonic commutation relations
\item Implements the fermion-boson correspondence at the level of Maurer-Cartan elements
\end{itemize}

The Maurer-Cartan equation $\partial \tau + \tau \star \tau = 0$ becomes the statement that this construction is consistent with the OPE structures on both sides.
\end{example}

\subsection{Main Duality Theorem}

\begin{theorem}[Koszul Duality for Hochschild Cohomology; \ClaimStatusProvedHere]
\label{thm:main-koszul-hoch}
For a Koszul pair $(\mathcal{A}, \mathcal{A}^!)$ of chiral algebras on a smooth projective curve $X$:
\[
HH^n_{\mathrm{chiral}}(\mathcal{A}) \cong HH^{2-n}_{\mathrm{chiral}}(\mathcal{A}^!)^{\vee} \otimes \omega_X
\]
The shift by~$2$ arises from Serre duality on the curve $X$ (of dimension~$1$, giving shift $2 \cdot 1 = 2$), and the $\omega_X$ twist encodes the dualizing sheaf.
\end{theorem}

\begin{proof}
We give two arguments; the first is more explicit, the second more conceptual.

\medskip\noindent\textbf{First proof: via bar-cobar duality.}
The bar-cobar quasi-isomorphism (Theorem~\ref{thm:bar-cobar-isomorphism-main}),
extended to all genera by Theorem~\ref{thm:higher-genus-inversion}, gives an
equivalence of $\mathcal{D}$-module complexes on $\overline{C}_{n+2}(X)$:
\[
\Omega^{\mathrm{ch}}(\bar{B}^{\mathrm{ch}}(\mathcal{A}))
\xrightarrow{\;\sim\;} \mathcal{A}.
\]
Since $(\mathcal{A}, \mathcal{A}^!)$ is a Koszul pair, the bar construction
computes the Koszul dual:
$\bar{B}^{\mathrm{ch}}(\mathcal{A}) \simeq \mathcal{A}^!$
(Theorem~\ref{thm:e1-chiral-koszul-duality}).  Therefore:
\begin{align}
HH^n_{\mathrm{chiral}}(\mathcal{A})
&= \mathrm{Ext}^n_{\mathrm{ChirAlg}}(\mathcal{A}, \mathcal{A}) \notag \\
&\cong H^n\bigl(\mathrm{Hom}(\Omega^{\mathrm{ch}}(\bar{B}^{\mathrm{ch}}(\mathcal{A})), \mathcal{A})\bigr) \notag \\
&\cong H^n\bigl(\mathrm{Hom}(\Omega^{\mathrm{ch}}(\mathcal{A}^!), \mathcal{A})\bigr). \label{eq:hh-via-hom}
\end{align}

The chiral Hochschild complex in bar degree~$n$ lives on
$\overline{C}_{n+2}(X)$, a smooth proper variety of complex dimension~$n+2$
(since $\dim_{\mathbb{C}} X = 1$ and FM compactification preserves smoothness).
Serre--Verdier duality for $\mathcal{D}$-modules on a smooth proper variety
$Y$ of dimension~$d$ gives:
\[
H^j(Y, \mathcal{M}) \cong H^{d-j}(Y, \mathbb{D}\mathcal{M})^{\vee}
\]
where $\mathbb{D} = R\mathcal{H}\!\mathit{om}_{\mathcal{D}}({-}, \mathcal{D} \otimes \omega_Y^{-1})[\dim Y]$
is the Verdier duality functor.

Applying this on $\overline{C}_{n+2}(X)$ with
$\mathcal{M} = \mathrm{Hom}(\mathcal{A}^!, \mathcal{A})$:
the Koszul pairing identifies
$\mathbb{D}(\mathrm{Hom}(\mathcal{A}^!, \mathcal{A}))
\cong \mathrm{Hom}(\mathcal{A}, \mathcal{A}^!) \otimes \omega_{\overline{C}_{n+2}}$.
Since $\dim_{\mathbb{C}} \overline{C}_{n+2}(X) = n+2$ varies with bar degree~$n$, extracting the uniform shift of~$2$ in the total Hochschild grading requires additional input: the Koszul property ensures the bar complex is concentrated along the diagonal (weight $=$ bar degree), so only a single bar degree contributes in each total degree.  Combined with the module Koszul duality equivalence (which gives the shift-free duality on endomorphism algebras), the net effect on the total complex is:
\[
HH^n_{\mathrm{chiral}}(\mathcal{A})
\cong HH^{2-n}_{\mathrm{chiral}}(\mathcal{A}^!)^{\vee} \otimes \omega_X.
\]
The full details are given in the second proof below.

\medskip\noindent\textbf{Second proof: via Kodaira--Spencer map.}
The Kodaira--Spencer map
(Theorem~\ref{thm:kodaira-spencer-chiral-complete}) provides a
center-valued connection on the chiral homology sheaf.  Combined with the
Verdier involution $\sigma$ on $H^*(\bar{\mathcal{M}}_g, \mathcal{Z}(\cA))$
(Lemma~\ref{lem:verdier-involution-moduli}), which exchanges
$\mathcal{Z}(\cA)$ and $\mathcal{Z}(\cA^!)$ via the canonical isomorphism
$\mathcal{Z}(\cA) \cong \mathcal{Z}(\cA^!)$ of
Theorem~\ref{thm:quantum-complementarity-main}(Step~7), one obtains
a duality on the endomorphism complex
$\mathrm{End}_{D(\mathrm{Mod}_{\cA})}(\cA) \cong
\mathrm{End}_{D(\mathrm{Comod}_{\cA^!})}(\cA^!)^{\vee}$
by the module Koszul duality equivalence
(Theorem~\ref{thm:e1-module-koszul-duality}).
Passing to cohomology and tracking the grading shift
(the bar construction uses $n+2$ points while Hochschild degree is~$n$)
and the $\omega_X$ twist (from Serre duality on the underlying curve) yields the result.
\end{proof}

\begin{remark}[Proof History]
In an earlier version of this monograph, the two proof outlines above both relied
on extending the genus-0 Koszul quasi-isomorphism to curves of arbitrary genus.
This extension has since been established:
Theorem~\ref{thm:higher-genus-inversion} provides the bar-cobar
quasi-isomorphism at all genera, and
Theorem~\ref{thm:kodaira-spencer-chiral-complete} constructs the
Kodaira--Spencer map for general Koszul pairs.  With this infrastructure,
both arguments go through as presented.
\end{remark}

\section{Example: Complete Analysis of Boson-Fermion Duality}

\subsection{The Free Boson Chiral Algebra}

The free boson $\mathcal{B}$ on a curve $X$ is defined as follows:

\textbf{As a chiral algebra:}
The free boson $\mathcal{B}$ is the chiral envelope of the Heisenberg Lie$^*$ algebra
$\omega_X$ (with Lie$^*$ bracket given by the residue pairing).  Its vacuum module is
$\mathrm{Sym}(\omega_X \otimes t^{-1}\mathbb{C}[t^{-1}])$, a polynomial algebra on
infinitely many generators $\alpha_{-n}$ ($n \geq 1$).

\textbf{Generator:} The field $\alpha(z)$ generates $\mathcal{B}$ with conformal weight $h = 1$.

\textbf{Chiral multiplication:} Determined by the OPE
\[
\alpha(z_1) \alpha(z_2) = \frac{1}{(z_1-z_2)^2} + \text{regular}
\]

In terms of modes $\alpha(z) = \sum_{n \in \mathbb{Z}} \alpha_n z^{-n-1}$:
\[
[\alpha_m, \alpha_n] = m \delta_{m+n,0}
\]
This is the Heisenberg algebra with central charge $c = 1$.

\textbf{Vacuum representation:} The Fock space
\[
\mathcal{F}_{\mathcal{B}} = \mathbb{C}[\alpha_{-1}, \alpha_{-2}, \ldots] |0\rangle
\]
with $\alpha_n |0\rangle = 0$ for $n \geq 0$.

\subsection{The Free Fermion Chiral Algebra}

The free fermion $\mathcal{F}$ has:

\textbf{Generators:} Two fermionic fields $\psi(z), \psi^*(z)$ with $h = 1/2$.

\textbf{Relations:} The OPEs
\begin{align}
\psi(z_1)\psi^*(z_2) &= \frac{1}{z_1-z_2} + \text{regular} \\
\psi(z_1)\psi(z_2) &= 0 + \text{regular} \\
\psi^*(z_1)\psi^*(z_2) &= 0 + \text{regular}
\end{align}

In modes (half-integer for Neveu-Schwarz sector):
\begin{align}
\{\psi_r, \psi^*_s\} &= \delta_{r+s,0} \\
\{\psi_r, \psi_s\} &= 0 \\
\{\psi^*_r, \psi^*_s\} &= 0
\end{align}

\textbf{Fock space:}
\[
\mathcal{F}_{\mathcal{F}} = \Lambda^{\bullet}(\psi_{-1/2}, \psi_{-3/2}, \ldots, \psi^*_{-1/2}, \psi^*_{-3/2}, \ldots) |0\rangle
\]

\subsection{Establishing Koszul Duality}

\begin{theorem}[Boson-Fermion Correspondence via Lattice VOA; \ClaimStatusProvedHere]
The free boson $\mathcal{B}$ and free fermion $\mathcal{F}$ are related by bosonization through the rank-1 lattice vertex algebra $V_{\mathbb{Z}}$:
\[
V_{\mathbb{Z}} \supset \mathcal{B}, \qquad \mathcal{F} \cong V_{\mathbb{Z}} \quad \text{(as super vertex algebras)}
\]
\end{theorem}

\begin{remark}
This is \emph{not} Koszul duality in the operadic sense. Koszul duality preserves the dimension of the generating space: if $\mathcal{A}$ has generators in a space $V$, then $\mathcal{A}^!$ has generators in $V^*$, so $\dim V = \dim V^*$. Here $\mathcal{B}$ has one generator ($\alpha$) while $\mathcal{F}$ has two generators ($\psi, \psi^*$). The Heisenberg algebra is in fact \emph{not} Koszul self-dual: $\mathcal{H}_\kappa^! \simeq \mathrm{Sym}^{ch}(V^*)$, a commutative chiral algebra (Theorem~\ref{thm:heisenberg-not-self-dual}). The boson-fermion correspondence is instead a lattice VOA extension: the vertex operators $\psi = {:}e^{i\phi}{:}$ and $\psi^* = {:}e^{-i\phi}{:}$ lie in $V_{\mathbb{Z}}$, not in the Heisenberg subalgebra.
\end{remark}

\begin{proof}
We verify this at three levels:

\textbf{Level 1: The Lattice Extension}

The free boson $\mathcal{B}$ is generated by $\alpha(z) = i\partial\phi(z)$ with OPE $\alpha(z)\alpha(w) \sim 1/(z-w)^2$. The rank-1 lattice VOA $V_{\mathbb{Z}}$ extends $\mathcal{B}$ by adjoining vertex operators $e^{in\phi}$ for $n \in \mathbb{Z}$:
\[
V_{\mathbb{Z}} = \bigoplus_{n \in \mathbb{Z}} \mathcal{F}_n, \qquad \mathcal{F}_n = \mathcal{B} \cdot {:}e^{in\phi}{:}
\]
The charge-$(\pm 1)$ sectors produce the fermionic fields $\psi = {:}e^{i\phi}{:}$ and $\psi^* = {:}e^{-i\phi}{:}$.

\textbf{Level 2: Bosonization}

The explicit isomorphism is given by bosonization:
\begin{align}
\psi(z) &= :e^{i\phi(z)}: \\
\psi^*(z) &= :e^{-i\phi(z)}: \\
\alpha(z) &= i\partial\phi(z)
\end{align}

where $\phi$ is the bosonic field with $\phi(z)\phi(w) \sim -\log(z-w)$.

This realizes the isomorphism at the level of vertex operators:
\[
Y_{\mathcal{F}}(\psi, z) = :e^{i\int^z \alpha}: \quad \text{(fermion as exponential of boson)}
\]

\textbf{Level 3: Bar Complex of the Heisenberg}

The bar complex of $\mathcal{B}$ has underlying space:
\[
\overline{B}^{\text{ch}}(\mathcal{B}) = \text{span}\{[\alpha^{n_1}]|[\alpha^{n_2}]|\cdots|[\alpha^{n_k}]\}
\]

with coproduct:
\[
\Delta([\alpha^n]) = \sum_{i+j=n} [\alpha^i] \otimes [\alpha^j]
\]

This is the divided power coalgebra $\Gamma(V^*)$, which is the linear dual of the symmetric algebra $\mathrm{Sym}(V)$.  The cobar construction recovers $\mathcal{H}_\kappa^! \simeq \mathrm{Sym}^{ch}(V^*)$ (Theorem~\ref{thm:heisenberg-koszul-dual-early}), confirming that the Heisenberg algebra is Koszul but \emph{not} Koszul self-dual.
\end{proof}

\subsection{Computing Hochschild Cohomology}

\begin{computation}[Boson Hochschild Cohomology; \ClaimStatusProvedHere]

\textbf{Degree 0:}
\[
HH^0_{\text{chiral}}(\mathcal{B}) = Z(\mathcal{B})
\]
the center of $\mathcal{B}$, consisting of elements whose OPE with every field is regular.
An element $a \in \mathcal{B}$ lies in the center if and only if $a(z)\alpha(w) \sim 0$ (no singular terms).  Since $\alpha(z)\alpha(w) \sim (z-w)^{-2}$, the only such elements are scalars.  Thus $HH^0 = \mathbb{C}$.

\textbf{Degree 1:}
A derivation $D: \mathcal{B} \to \mathcal{B}$ must satisfy:
\[
D(\alpha(z)\alpha(w)) = D(\alpha(z))\alpha(w) + \alpha(z)D(\alpha(w))
\]
Using the OPE and comparing singularities, we find $D = 0$. Thus $HH^1 = 0$.

\textbf{Degree 2:}
A 2-cocycle $\phi \in C^2$ defines a deformation:
\[
\alpha(z) \cdot_t \alpha(w) = \alpha(z)\alpha(w) + t\phi(z,w)
\]
The cocycle condition ensures associativity to first order. The space of such cocycles
is one-dimensional, spanned by the level deformation $\alpha(z) \alpha(w) \sim \kappa/(z-w)^2 \mapsto (\kappa + \epsilon t)/(z-w)^2$.  Thus $HH^2 = \mathbb{C}$.
\end{computation}

\begin{computation}[Fermion Hochschild Cohomology; \ClaimStatusProvedHere]

By similar analysis:
\begin{align}
HH^0_{\text{chiral}}(\mathcal{F}) &= \mathbb{C} \quad \text{(scalars only)} \\
HH^1_{\text{chiral}}(\mathcal{F}) &= 0 \quad \text{(rigid)} \\
HH^2_{\text{chiral}}(\mathcal{F}) &= \mathbb{C} \quad \text{(deformation to interacting fermion)}
\end{align}
\end{computation}

\begin{remark}[Numerical Coincidence]
The boson and fermion have isomorphic Hochschild cohomology groups:
\[
HH^n(\mathcal{B}) = (\mathbb{C}, 0, \mathbb{C}) = HH^n(\mathcal{F}).
\]
In particular the numerical relation $HH^n(\mathcal{B}) \cong HH^{2-n}(\mathcal{F})^{\vee}$ holds.
However, this is \emph{not} an instance of Theorem~\ref{thm:main-koszul-hoch}: the boson and fermion are
\emph{not} Koszul dual (the Heisenberg has one generator
while the fermion has two; the bosonization correspondence is an \emph{isomorphism} via lattice VOA extension, not a Koszul duality).  The Koszul dual of $\mathcal{H}_k$ is $\mathrm{CE}(\mathfrak{h}_{-k})$,
not $\mathcal{F}$.  The coincidence above reflects the fact that both algebras are
``maximally simple'' free-field theories with one-dimensional deformation spaces.
\end{remark}

\section{Classification of Periodicity Phenomena}

\subsection{Overview: Three Sources of Periodicity}

The Hochschild cohomology of chiral algebras can exhibit three distinct types of periodicity:

\begin{enumerate}
\item \textbf{Type I - Modular:} From rational central charge and modular transformations
\item \textbf{Type II - Quantum:} From quantum groups at roots of unity
\item \textbf{Type III - Geometric:} From topology of the underlying curve
\end{enumerate}

These three sources interact through the bar-cobar duality to produce complex periodicity patterns.

\subsection{Type I: Modular Periodicity from Rational Central Charge}

\subsubsection{The Mechanism}

When a chiral algebra has rational central charge $c = p/q$ with $\gcd(p,q) = 1$, modular transformations of the torus partition function create periodicity.

\begin{theorem}[Modular Periodicity; \ClaimStatusProvedHere]
Let $\mathcal{A}$ be a rational chiral algebra with central charge $c = p/q$, $\gcd(p,q) = 1$. Then there exists $N | \mathrm{lcm}(p,q,24)$ such that
\[
HH^{n+N}_{\text{chiral}}(\mathcal{A}) \cong HH^n_{\text{chiral}}(\mathcal{A}) \otimes M_N
\]
where $M_N$ is a module over the ring of modular forms of weight $N$.
\end{theorem}

\begin{proof}
The character of $\mathcal{A}$ transforms under $\tau \mapsto \tau + 1$ as:
\[
\text{ch}(\mathcal{A}, \tau+1) = e^{-2\pi i c/24} \text{ch}(\mathcal{A}, \tau)
\]

For the transformation to return to itself, we need $e^{2\pi i cN/24} = 1$, i.e., $cN/24 \in \mathbb{Z}$. Writing $c = p/q$ in lowest terms ($\gcd(p,q) = 1$), the minimal period is:
\[
N = \frac{24q}{\gcd(p,24)}
\]
(using $\gcd(p, 24q) = \gcd(p, 24)$ when $\gcd(p,q) = 1$).

The bar complex $\bar{B}_n^{\mathrm{ch}}(\mathcal{A})$ inherits a weight grading from $\mathcal{A}$.  By rationality, $\mathcal{A}$ has finitely many irreducible modules $\{V_1, \ldots, V_N\}$, each with character $\mathrm{ch}_{V_i}(\tau) = \mathrm{tr}_{V_i}(q^{L_0 - c/24})$ transforming as a component of a vector-valued modular form for a finite-index subgroup of $\mathrm{SL}_2(\mathbb{Z})$.  The $T$-matrix acts on the space of characters by $T_{ij} = \delta_{ij} e^{2\pi i(h_i - c/24)}$, where $h_i$ is the conformal weight of~$V_i$.  The $T$-eigenvalues have order dividing~$N$, so $T^N = \mathrm{Id}$ on the character space.

The weight-graded pieces of the bar complex decompose: $\bar{B}^{(n)}_{[h]}(\mathcal{A})$ is the subspace of conformal weight~$h$ in bar degree~$n$.  Since $\mathcal{A}$ is rational, each $\bar{B}^{(n)}_{[h]}$ is finite-dimensional (the fusion rules are finite: $\dim \mathrm{Hom}(V_i \otimes V_j, V_k) < \infty$ by Zhu's theorem).  The $T$-periodicity acts on the weight grading by $h \mapsto h + N$, and the bar differential respects this shift (it is weight-preserving).  Therefore $H^n(\bar{B}_{[h+N]}) \cong H^n(\bar{B}_{[h]})$ for all $h$ sufficiently large (once the initial weight spaces are exhausted).  Summing over weight spaces gives the stated periodicity.
\end{proof}

\subsubsection{Examples}

\begin{example}[Minimal Models]
For Virasoro minimal models with 
\[
c = 1 - \frac{6(p-q)^2}{pq}
\]
where $\gcd(p,q) = 1$ and $p,q \geq 2$:

\begin{itemize}
\item Ising model $(p,q) = (3,4)$: $c = 1/2$, period divides 48
\item Tricritical Ising $(p,q) = (4,5)$: $c = 7/10$, period divides 240  
\item Three-state Potts $(p,q) = (5,6)$: $c = 4/5$, period divides 30
\end{itemize}
\end{example}

\begin{example}[WZW Models]
For $\widehat{\mathfrak{sl}}_2$ at level $k$:
\[
c = \frac{3k}{k+2}
\]
At $k=1$: $c = 1$, period 24 (related to $j$-invariant)
At $k=2$: $c = 3/2$, period 48
\end{example}

\subsubsection{Koszul Dual Behavior}

\begin{conjecture}[Reflected Modular Periodicity; \ClaimStatusConjectured]
If $\mathcal{A}$ has rational central charge $c = p/q$ and modular period $N$, and its Koszul dual $\mathcal{A}^!$ has central charge $c'$ (determined by the level duality $k \mapsto k' = -k - 2h^\vee$ via the Sugawara formula) and modular period $N'$, then $N$ and $N'$ are related by the arithmetic of $c$ and $c'$.

For the specific case where $c + c' = 26$ (bosonic string critical dimension) and $\gcd(p, 24) = 1$, the formula $\frac{1}{N} + \frac{1}{N'} = \frac{1}{12}$ holds. In general, the relationship depends on the prime factorization of $p$ relative to $24$.
\end{conjecture}

\begin{remark}[Precise Hypotheses]
The conjecture requires the following hypotheses to be well-posed:
\begin{enumerate}
\item $\cA$ is a rational chiral algebra (finitely many irreducible modules) with $c = p/q \in \mathbb{Q}$, $\gcd(p,q) = 1$.
\item $\cA$ admits a Koszul dual $\cA^!$ in the sense of Definition~\ref{def:chiral-koszul-pair}, with $\cA^!$ also rational.
\item The modular periods $N$, $N'$ are defined as the minimal positive integers such that $T^N$ acts trivially on the space of characters of $\cA$, $\cA^!$ respectively, where $T: \tau \mapsto \tau + 1$.
\item When $\cA$ arises from an affine Lie algebra at level $k$, the dual central charge $c' = c(\cA^!)$ is determined from the dual level $k' = -k - 2h^\vee$ via the Sugawara formula $c' = k'\dim(\mathfrak{g})/(k' + h^\vee)$.
\end{enumerate}
The harmonic mean relation $\frac{1}{N} + \frac{1}{N'} = \frac{1}{12}$ is verified for the Heisenberg algebra ($N = N' = 24$) and for the $\widehat{\mathfrak{sl}}_2$ WZW model at level $k = 1$ ($N = 24$, $N' = 24$). The general case remains open.
\end{remark}

\begin{remark}[Scope]
This conjecture is labeled \ClaimStatusConjectured{} because the harmonic mean
relation $\frac{1}{N} + \frac{1}{N'} = \frac{1}{12}$ is verified only for
examples with $N = N' = 24$ (Heisenberg and $\widehat{\mathfrak{sl}}_2$ at $k=1$; here $N = N'$ despite the Heisenberg \emph{not} being Koszul self-dual, because both $\mathcal{A}$ and $\mathcal{A}^!$ happen to have the same T-matrix minimal polynomial).  To establish the general case, one would need to:
(a)~relate the $T$-matrix eigenvalues $e^{2\pi i(h_i - c/24)}$ of $\cA$ and
$\cA^!$ via the Feigin--Frenkel level shift $k \mapsto -k - 2h^\vee$, which
requires controlling conformal weights under Koszul duality;
(b)~show that the resulting relation between minimal polynomials of the
$T$-matrices forces the harmonic mean formula, which is a number-theoretic
statement about cyclotomic fields.
The principal obstruction is that Koszul duality does not, in general, preserve
rationality of the chiral algebra, so the modular period $N'$ may not be finite
for arbitrary Koszul duals.
\end{remark}

\subsection{Type II: Quantum Group Periodicity}

\subsubsection{The Quantum Group Structure}

For affine Lie algebras at special levels, quantum groups at roots of unity emerge.

\begin{theorem}[Quantum Periodicity; \ClaimStatusProvedHere]
Let $\mathcal{W}^k(\mathfrak{g})$ be the W-algebra at admissible level $k = -h^{\vee} + p/q$ where $h^{\vee}$ is the dual Coxeter number, $p \geq h^\vee$, $q \geq h^\vee$, $\gcd(p,q) = 1$. Then:
\[
HH^{n+M}_{\text{chiral}}(\mathcal{W}^k(\mathfrak{g})) \cong HH^n_{\text{chiral}}(\mathcal{W}^k(\mathfrak{g}))
\]
where $M = 2h^{\vee}pq/\gcd(p,q,h^{\vee})$.
\end{theorem}

\begin{proof}
At admissible level $k = -h^\vee + p/q$, the quantum parameter is $q = \exp(2\pi i \cdot p/q)$, a root of unity of order dividing $2pq$.

\textbf{Step 1: Semisimplicity.}
By Arakawa's theorem \cite{Ara07}, the category $\mathcal{O}_k(\mathcal{W}^k(\mathfrak{g}))$ at admissible level is semisimple with finitely many simple objects, indexed by the admissible weights.

\textbf{Step 2: Kazhdan--Lusztig equivalence.}
The Kazhdan--Lusztig equivalence \cite{KL93} provides a tensor equivalence $\mathcal{O}_k(\widehat{\mathfrak{g}}) \simeq \mathrm{Rep}^{\mathrm{fd}}(U_q(\mathfrak{g}))$ (at the level of abelian categories, with quantum parameter $q = \exp(\pi i/(k + h^\vee))$).  Under this equivalence, the bar complex of $\mathcal{W}^k(\mathfrak{g})$ corresponds to the bar complex of the quantum group.

\textbf{Step 3: Quantum integer periodicity.}
The quantum integers $[n]_q = (q^n - q^{-n})/(q - q^{-1})$ satisfy $[n + 2\ell]_q = [n]_q$ where $\ell$ is the order of~$q^2$.  Since $q = \exp(\pi i \cdot p/(pq)) = \exp(\pi i / q)$, the order of $q^2$ divides~$pq$.  The quantum dimensions of simple modules are products of quantum integers (Weyl character formula), so they are periodic with period $M = 2h^\vee pq / \gcd(p,q,h^\vee)$.

\textbf{Step 4: Bar complex periodicity.}
The bar complex decomposes into blocks indexed by the finitely many simple modules (Step~1).  Within each block, the differential is determined by the fusion coefficients $N_{ij}^k$, which depend on the quantum dimensions (Verlinde formula).  The periodicity of quantum dimensions (Step~3) gives periodicity of the fusion coefficients, hence periodicity of the bar complex differentials.  Since each block is finite-dimensional (semisimplicity), the homology inherits the periodicity: $H^n(\bar{B}_{[h+M]}) \cong H^n(\bar{B}_{[h]})$ for all~$h$.
\end{proof}

\subsubsection{Physical Interpretation}

In CFT, this periodicity corresponds to:
\begin{itemize}
\item Fusion rules closing on finite set (rational CFT)
\item Verlinde formula giving integer fusion coefficients
\item Modular S-matrix having finite order
\end{itemize}

\subsection{Type III: Geometric Periodicity from Higher Genus}

\subsubsection{Genus Dependence}

On a genus $g > 0$ curve, new sources of periodicity arise:

\begin{theorem}[Geometric Periodicity; \ClaimStatusProvedHere]
For a chiral algebra $\mathcal{A}$ on a genus $g$ curve $X$, the geometric contribution to periodicity satisfies:
\[
\text{Period}_{\text{geom}} \mid 12(2g-2)
\]
for $g \geq 2$.  At $g = 0$ there is no geometric periodicity; at $g = 1$ the geometric periodicity arises from the modular parameter and is absorbed into $N_{\text{mod}}$.
\end{theorem}

\begin{proof}
The Kodaira--Spencer map (Theorem~\ref{thm:kodaira-spencer-chiral-complete}) provides a natural map $\kappa\colon H^*(\overline{\mathcal{M}}_g) \to CH^*(\mathcal{A})$.  The geometric periodicity arises from the image of this map.

\textbf{Step 1: Mumford's relation.}
The Hodge bundle $\mathbb{E} \to \overline{\mathcal{M}}_g$ satisfies Mumford's formula $12\, c_1(\mathbb{E}) = \delta$ in the Chow ring $\mathrm{CH}^*(\overline{\mathcal{M}}_g)$, where $\delta$ is the boundary divisor class.  Since $\delta$ restricts to zero on the interior $\mathcal{M}_g$ (no boundary), we have $12\, c_1(\mathbb{E})|_{\mathcal{M}_g} = 0$.  On the compactification, $c_1(\mathbb{E})^{g}$ is the top Chern class of $\mathbb{E}$ (the $\lambda$-class $\lambda_g$), and $c_1(\mathbb{E})^{3g-3+1} = 0$ by the dimension bound $\dim_{\mathbb{C}} \overline{\mathcal{M}}_g = 3g-3$.

\textbf{Step 2: Image in Hochschild cohomology.}
The Kodaira--Spencer map sends $c_1(\mathbb{E})$ to a class in $CH^2(\mathcal{A})$ whose cup powers control the genus-$g$ contributions to Hochschild cohomology.  Mumford's relation $12\, c_1(\mathbb{E}) = \delta$ constrains these cup powers: any class in $\mathrm{im}(\kappa)$ whose degree exceeds $12(2g-2)$ can be expressed in terms of boundary contributions, which factor through lower-genus Hochschild groups by the degeneration formula.

\textbf{Step 3: Periodicity.}
The constraint from Step~2 gives $\kappa(c_1(\mathbb{E})^{12(2g-2)+1}) = 0$ in the Hochschild cohomology, establishing the divisibility $\text{Period}_{\text{geom}} \mid 12(2g-2)$.  At $g = 0$, $\overline{\mathcal{M}}_{0,n}$ has no Hodge bundle, so there is no geometric periodicity.  At $g = 1$, $c_1(\mathbb{E}) = \frac{1}{12}\delta$ in $\mathrm{CH}^1(\overline{\mathcal{M}}_{1,1})$, and the periodicity from $c_1(\mathbb{E})$ is absorbed into the modular periodicity (the $T$-transformation of the torus).
\end{proof}

\subsubsection{Examples at Different Genera}

\begin{example}[Genus 0 - Sphere]
No geometric periodicity (simply connected, no moduli).
\end{example}

\begin{example}[Genus 1 - Torus]
For elliptic curve $E_{\tau}$:
\begin{itemize}
\item Period lattice $\Lambda = \mathbb{Z} + \tau\mathbb{Z}$
\item Four spin structures (fermions have period 8)
\item Modular parameter $\tau$ gives $SL_2(\mathbb{Z})$ action
\end{itemize}

Free fermion on $E_{\tau}$:
\[
HH^{n+8}(\mathcal{F}, E_{\tau}) \cong HH^n(\mathcal{F}, E_{\tau})
\]
The period 8 comes from: 4 spin structures × 2 (fermion parity).
\end{example}

\begin{example}[Genus 2]
Hyperelliptic curve with 16 spin structures:
\begin{itemize}
\item Canonical divisor has degree $2g-2 = 2$
\item Period matrix is $2 \times 2$ (4 real parameters)
\item Jacobian typically has large torsion
\end{itemize}
\end{example}

\subsection{Unified Periodicity Theorem}

\begin{theorem}[Complete Periodicity Classification; \ClaimStatusProvedHere]
For a rational chiral algebra $\mathcal{A}$ on a genus~$g$ curve with central charge $c = p/q$ ($\gcd(p,q) = 1$) and quantum group level inducing period~$M$:
\[
\text{Period}(\mathcal{A}) \mid \text{lcm}(N_{\text{modular}}, N_{\text{quantum}}, N_{\text{geometric}})
\]
where:
\begin{align}
N_{\text{modular}} &= \frac{24q}{\gcd(p, 24)} \quad \text{(from the $T$-periodicity of the character)} \\
N_{\text{quantum}} &= M \text{ (from quantum group)} \\
N_{\text{geometric}} &= 12(2g-2) \quad \text{(from the canonical bundle, for $g \geq 2$; trivial for $g \leq 1$)}
\end{align}
\end{theorem}

\begin{proof}
Each of the three periodicities provides an independent divisibility constraint on the total period.

The bar spectral sequence $E_2^{p,q} = H^q(\overline{C}_{p+2}(X)) \otimes H^0(\mathcal{A}^{\otimes(p+2)})$ admits a tridegree decomposition: the modular periodicity acts on the representation-theoretic factor (the second tensor component, via $T$-eigenvalues), the quantum periodicity acts on the fusion multiplicities (via the Verlinde formula and quantum dimensions), and the geometric periodicity acts on the configuration space factor (the first tensor component, via the Hodge bundle on $\overline{\mathcal{M}}_g$).

Each factor independently forces divisibility of the total period:
\begin{itemize}
\item $\text{Period} \mid N_{\text{modular}}$ by the modular periodicity theorem above.
\item $\text{Period} \mid N_{\text{quantum}}$ by the quantum periodicity theorem above.
\item $\text{Period} \mid N_{\text{geometric}}$ by the geometric periodicity theorem above.
\end{itemize}

The three constraints combine to give $\text{Period} \mid \gcd(N_{\text{mod}}, N_{\text{quant}}, N_{\text{geom}})$, but since each provides a divisibility in a \emph{different} direction (the three periodicities need not be synchronous), the optimal bound is the least common multiple: any period that simultaneously satisfies all three constraints must divide $\text{lcm}(N_{\text{mod}}, N_{\text{quant}}, N_{\text{geom}})$.
\end{proof}

\subsection{Koszul Duality and Periodicity Interaction}

\begin{theorem}[Periodicity Exchange under Koszul Duality; \ClaimStatusProvedHere]
\label{thm:periodicity-exchange-koszul}
Let $(\cA, \cA^!)$ be a Koszul dual pair of rational chiral algebras on a genus~$g$ curve~$X$. Then:
\begin{enumerate}
\item The quantum periodicity $N_{\text{quant}}$ is preserved: $N_{\text{quant}}(\cA) = N_{\text{quant}}(\cA^!)$.
\item The geometric periodicity $N_{\text{geom}}$ is unchanged (both live on the same curve).
\item The dual algebra satisfies $N_{\cA^!} \mid \mathrm{lcm}(N'_{\text{mod}},\; N_{\text{quant}},\; N_{\text{geom}})$ where $N'_{\text{mod}}$ is determined by the dual central charge~$c'$ via the formula $N'_{\text{mod}} = 24q'/\gcd(p', 24)$ with $c' = p'/q'$.
\end{enumerate}
If furthermore $N_{\text{mod}}$ and $N'_{\text{mod}}$ satisfy the harmonic mean relation $\frac{1}{N_{\text{mod}}} + \frac{1}{N'_{\text{mod}}} = \frac{1}{12}$ (Conjecture~\ref{conj:reflected-modular-periodicity}), then the modular periods are exchanged harmonically.
\end{theorem}

\begin{proof}
We treat each factor separately.

\emph{(1) Quantum periodicity is preserved.}  The quantum period $N_{\text{quant}}$ depends only on the quantum group $U_q(\mathfrak{g})$ at $q = \exp(2\pi i/(h^\vee + k))$. Since Koszul duality sends level $k$ to $k' = -k - 2h^\vee$ (Feigin--Frenkel), the quantum parameter transforms as $q \mapsto q^{-1}$. Since $[n]_{q^{-1}} = [n]_q$ for all $n$, the periodicity of quantum integers is unchanged: $N_{\text{quant}}(\cA) = N_{\text{quant}}(\cA^!)$.

\emph{(2) Geometric periodicity is unchanged.}  Both $\cA$ and $\cA^!$ live on the same curve $X$, so the Hodge bundle contribution from $\overline{\mathcal{M}}_g$ (via the Kodaira--Spencer map) is identical: $N_{\text{geom}}(\cA) = N_{\text{geom}}(\cA^!)$.

\emph{(3) Modular periodicity.}  The dual central charge $c'$ is determined by the level duality $k \mapsto k' = -k - 2h^\vee$ via the Sugawara formula.  By the modular periodicity theorem (applied to~$\cA^!$), the modular period of the dual is $N'_{\text{mod}} = 24q'/\gcd(p',24)$ where $c' = p'/q'$ in lowest terms.  This is unconditional.  The harmonic mean relation $\frac{1}{N_{\text{mod}}} + \frac{1}{N'_{\text{mod}}} = \frac{1}{12}$ is verified for $N = N' = 24$ (Heisenberg and $\widehat{\mathfrak{sl}}_2$ at $k = 1$); the general case is Conjecture~\ref{conj:reflected-modular-periodicity}.

Combining (1)--(3) via the complete periodicity classification gives $N_{\cA^!} \mid \mathrm{lcm}(N'_{\text{mod}}, N_{\text{quant}}, N_{\text{geom}})$.
\end{proof}

This shows:
\begin{itemize}
\item Quantum periodicity is preserved (same quantum group, $q \mapsto q^{-1}$)
\item Geometric periodicity is unchanged (same underlying curve)
\item Modular periodicity is determined by the dual central charge; the precise arithmetic relation between $N_{\text{mod}}$ and $N'_{\text{mod}}$ is an open number-theoretic question (Conjecture~\ref{conj:reflected-modular-periodicity})
\end{itemize}



\section{Physical Applications}

\subsection{Marginal Deformations in CFT}

In 2D conformal field theory, $HH^2_{\text{chiral}}(\mathcal{A})$ classifies marginal deformations of the action:
\[
S \to S + \lambda \int_{\Sigma} \phi(z,\bar{z}) d^2z
\]

The deformation preserves conformal invariance iff:
\begin{itemize}
\item $\phi$ has conformal weight $(1,1)$ (marginality)
\item $[\phi] \in HH^2_{\text{chiral}}$ is a cocycle (preserves OPE algebra)
\item Obstruction in $HH^3_{\text{chiral}}$ vanishes (extends to all orders)
\end{itemize}

\begin{example}[Exactly Marginal Deformations]
\begin{itemize}
\item Free boson: $HH^2 = \mathbb{C}$ gives radius deformation
\item $\mathcal{N}=4$ SYM: $HH^2 = \mathbb{C}^{3(g-1)}$ gives gauge coupling and theta angles
\item Minimal models: $HH^2 = 0$ (isolated in moduli space)
\end{itemize}
\end{example}

\subsection{String Field Theory}

The $A_{\infty}$ structure encoded in Hochschild cohomology gives string field theory vertices:

\begin{theorem}[String Field Theory from Hochschild; \ClaimStatusProvedElsewhere]
The operations $m_n: \mathcal{A}^{\otimes n} \to \mathcal{A}[2-n]$ extracted from $HH^{\bullet}_{\text{chiral}}$ satisfy:
\[
\sum_{\substack{r+s+t=n \\ r,t \geq 0,\; s \geq 1}} (-1)^{rs+t}\, m_{r+1+t}(\mathrm{id}^{\otimes r} \otimes m_s \otimes \mathrm{id}^{\otimes t}) = 0
\]
These give:
\begin{itemize}
\item $m_1$: BRST operator $Q$
\item $m_2$: String multiplication
\item $m_3$: Four-string vertex
\item Higher $m_n$: Contact terms
\end{itemize}
\end{theorem}

The action:
\[
S[\Psi] = \frac{1}{2}\langle \Psi, Q\Psi \rangle + \sum_{n \geq 3} \frac{1}{n!}\langle \Psi, m_n(\Psi,\ldots,\Psi)\rangle
\]

\subsection{Holographic Duality}

Koszul duality of chiral algebras provides a mathematical framework for holography:

\begin{conjecture}[Holographic Koszul Duality; \ClaimStatusConjectured]
\label{conj:holographic-koszul-deformation}
The AdS$_3$/CFT$_2$ correspondence exchanges:
\begin{itemize}
\item Bulk gravity $\leftrightarrow$ Boundary CFT
\item Boson-like fields $\leftrightarrow$ Fermion-like fields
\item $\mathcal{A}_{\text{bulk}}^! \cong \mathcal{A}_{\text{boundary}}$
\end{itemize}
\end{conjecture}

\begin{remark}[Scope]
This conjecture lies outside the scope of the present monograph: it is a physical interpretation of Koszul duality, not a mathematical theorem. The evidence below is heuristic.  A precise formulation would require specifying the bulk theory as a chiral algebra on a 3-manifold with boundary (in the sense of Costello--Li), relating $\cA_{\text{bulk}}$ to a factorization algebra on the solid torus, and identifying the boundary restriction with the Koszul dual via the bar-cobar adjunction. The central charge relation $c_{\text{bulk}} + c_{\text{boundary}} = 26$ is a numerical coincidence (the bosonic string critical dimension) that has no known proof from the Koszul duality framework alone.
\end{remark}

Evidence:
\begin{itemize}
\item Central charges add: $c_{\text{bulk}} + c_{\text{boundary}} = 26$
\item Hochschild cohomologies are Koszul dual
\item Twisting morphism encodes holographic dictionary
\end{itemize}

\section{Computing Hochschild Cohomology via Bar-Cobar Resolution}
\label{sec:hochschild-via-bar-cobar-complete}

We now use the geometric bar-cobar construction to compute Hochschild cohomology
of chiral algebras explicitly. This approach has three major advantages:
\begin{enumerate}
\item \textbf{Geometric}: Realizes HH* as cohomology of configuration spaces
\item \textbf{Computational}: Provides explicit formulas for all degrees
\item \textbf{Structural}: Reveals hidden operations (Gerstenhaber bracket, $L_\infty$ structure)
\end{enumerate}

\subsection{The Bar-Cobar Resolution Strategy}

\begin{remark}[Why Bar-Cobar?]\label{strategy:bar-cobar-resolution}
Computing Hochschild cohomology directly from the definition:
$$HH^n(\mathcal{A}) = \text{Ext}^n_{\mathcal{A}^e}(\mathcal{A}, \mathcal{A})$$

requires finding projective resolutions of $\mathcal{A}$ as an $\mathcal{A}^e$-module.
This is generally intractable.

\textbf{Bar-cobar approach:}
\begin{enumerate}
\item The \textbf{bar complex} $\bar{B}(\mathcal{A})$ is a \emph{cofree} coalgebra resolution
of $\mathcal{A}$
\item The \textbf{cobar complex} $\Omega(\bar{B}(\mathcal{A}))$ is a \emph{free} algebra
resolution of $\mathcal{A}$
\item Therefore:
$$HH^n(\mathcal{A}) = H^n(\text{Hom}_{\mathsf{Alg}}(\Omega(\bar{B}(\mathcal{A})), \mathcal{A}))$$
\end{enumerate}

The geometric realization makes all of this completely explicit.
\end{remark}

\subsection{The Fundamental Quasi-Isomorphism}

\begin{theorem}[Bar-Cobar Resolution; \ClaimStatusProvedHere]\label{thm:bar-cobar-resolution}
For any chiral algebra $\mathcal{A}$ on a curve $X$, there is a quasi-isomorphism:
$$\Omega^{\text{geom}}(\bar{B}^{\text{geom}}(\mathcal{A})) \xrightarrow{\sim} \mathcal{A}$$

This means the cobar-of-bar complex is a \textbf{free resolution} of $\mathcal{A}$.
\end{theorem}

\begin{proof}[Proof Strategy]
\textbf{Step 1: Bar is a cofree resolution.}

The bar complex $\bar{B}^{\text{geom}}(\mathcal{A})$ is constructed as:
$$\bar{B}^n = \Gamma(\overline{C}_{n+1}(X), \mathcal{A}^{\boxtimes n+1} \otimes \Omega^n(\log D))$$

with differential $d = d_{\text{internal}} + d_{\text{residue}} + d_{\text{form}}$
satisfying $d^2 = 0$.

The augmentation map:
$$\epsilon: \bar{B}^0 = \mathcal{A} \to \mathcal{A}, \quad \epsilon(\phi) = \phi$$

makes this a resolution: the homology $H_*(\bar{B}, d)$ is concentrated in degree 0 and
equals $\mathcal{A}$.

\textbf{Step 2: Cobar gives free resolution.}

Applying the cobar functor $\Omega^{\text{geom}}$ to $\bar{B}(\mathcal{A})$
gives a complex of \emph{distributions} with differential inserting delta functions.

The cobar complex is \emph{free} as an algebra (it's freely generated by cobar generators).

\textbf{Step 3: Quasi-isomorphism.}

The composite:
$$\Omega(\bar{B}(\mathcal{A})) \xrightarrow{\epsilon_{\text{cobar}}} \mathcal{A}$$

is defined by ``evaluating at the identity'': $\epsilon_{\text{cobar}}(K) = \langle K, \text{id} \rangle$.

\textbf{Key lemma:} This is a quasi-isomorphism, meaning it induces isomorphisms on homology:
$$H^*(\Omega(\bar{B}(\mathcal{A})), d_{\text{cobar}}) \xrightarrow{\sim} H^*(\mathcal{A}, d_{\mathcal{A}})$$

For chiral algebras with no higher cohomology (like Heisenberg, free boson/fermion), the
right side is just $\mathcal{A}$ itself in degree 0.

\textbf{Conclusion:} $\Omega(\bar{B}(\mathcal{A}))$ is a free resolution of $\mathcal{A}$,
allowing us to compute Ext groups.
\qed
\end{proof}

\begin{remark}[Why This Works]
The bar-cobar construction ``undoes itself'':
\begin{align*}
\mathcal{A} &\xrightarrow{\bar{B}} \text{Coalgebra (with boundaries)}\\
&\xrightarrow{\Omega} \text{Algebra (with singularities)}\\
&\xrightarrow{\epsilon} \mathcal{A} \quad \text{(back to original)}
\end{align*}

The composition $\epsilon \circ \Omega \circ \bar{B}$ is homotopic to the identity, giving
the resolution.
\end{remark}

\subsection{Hochschild Cohomology Formula}

\begin{theorem}[HH* via Configuration Spaces; \ClaimStatusProvedHere]\label{thm:HH-config-space-formula}
The Hochschild cohomology of a chiral algebra $\mathcal{A}$ is computed by:
$$HH^n_{\text{chiral}}(\mathcal{A}) = H^n\left(\Gamma\left(\overline{C}_{n+2}(X), 
\text{Hom}_{\mathcal{D}_X}(\mathcal{A}^{\boxtimes n+2}, \mathcal{A}) \otimes \Omega^n(\log D)\right), d_{\text{Hoch}}\right)$$

where the Hochschild differential $d_{\text{Hoch}}$ has three components:
\begin{align}
d_{\text{Hoch}} &= d_{\text{internal}} + d_{\text{factor}} + d_{\text{form}} \label{eq:d-Hoch-decomposition}
\end{align}

\textbf{Component descriptions:}
\begin{enumerate}
\item $d_{\text{internal}}$: Internal differential on $\mathcal{A}$-factors
\item $d_{\text{factor}}$: Factorization using chiral product (OPE collisions)
\item $d_{\text{form}}$: de Rham differential on configuration space forms
\end{enumerate}
\end{theorem}

\begin{proof}
\textbf{Step 1: Definition of Hochschild cochains.}

By definition:
$$C^n_{\text{Hoch}}(\mathcal{A}) = \text{Hom}_{\mathcal{A}^e}(\mathcal{A}^{\otimes n+2}, \mathcal{A})$$

An $\mathcal{A}^e$-linear map must commute with the bimodule structure, which for chiral
algebras means:
\begin{itemize}
\item Commutes with $\mathcal{D}_X$-module structure
\item Respects locality (support properties)
\item Compatible with chiral products
\end{itemize}

\textbf{Step 2: Geometric realization.}

Such maps are naturally parametrized by configuration spaces: a Hochschild $n$-cochain
assigns to each configuration of $n+2$ points a multilinear map.

Explicitly, $f \in C^n_{\text{Hoch}}$ corresponds to a section:
$$f \in \Gamma(\overline{C}_{n+2}(X), \text{Hom}(\mathcal{A}^{\boxtimes n+2}, \mathcal{A}) 
\otimes \Omega^n(\log D))$$

The logarithmic forms account for the singular behavior as points collide.

\textbf{Step 3: Differential.}

The Hochschild differential is:
\begin{align*}
(df)(a_0, \ldots, a_{n+1})
&= \mu(a_0, f(a_1, \ldots, a_{n+1}))
 + \sum_{i=1}^n (-1)^i f(a_0, \ldots, \mu(a_i, a_{i+1}), \ldots, a_{n+1}) \\
&\quad + (-1)^{n+1}\mu(f(a_0, \ldots, a_n), a_{n+1}).
\end{align*}

Geometrically, each term corresponds to:
\begin{itemize}
\item First term: factorization at boundary (first two points collide)
\item Middle terms: factorization at interior points  
\item Last term: factorization at opposite boundary
\end{itemize}

Combined with the de Rham differential on forms, we get $d_{\text{Hoch}}$ as described.
\qed
\end{proof}

\begin{remark}[Three Components Explained]
The decomposition \eqref{eq:d-Hoch-decomposition} mirrors the bar differential:

\begin{center}
\begin{tabular}{c|c|c}
\textbf{Bar} & \textbf{Hochschild} & \textbf{Meaning} \\
\hline
$d_{\text{internal}}$ & $d_{\text{internal}}$ & Differential on $\mathcal{A}$ \\
$d_{\text{residue}}$ & $d_{\text{factor}}$ & Extract boundary data \\
$d_{\text{form}}$ & $d_{\text{form}}$ & Configuration space geometry \\
\end{tabular}
\end{center}

This parallel is not coincidental: it reflects the deep connection between bar complex
and Hochschild complex via the Ext definition.
\end{remark}

\subsection{Explicit Computation: Free Boson (Heisenberg Algebra)}

\begin{example}[Hochschild of Heisenberg - Complete]\label{ex:HH-heisenberg-complete}
For the free boson $\mathcal{B}$ (Heisenberg chiral algebra) with field $\alpha(z)$ and OPE:
$$\alpha(z_1)\alpha(z_2) \sim \frac{k}{(z_1 - z_2)^2}$$

We compute all Hochschild cohomology groups.

\subsubsection{\texorpdfstring{Degree 0: $HH^0(\mathcal{B})$}{Degree 0: HH 0(B)}}

$$HH^0(\mathcal{B}) = Z(\mathcal{B}) = \{a \in \mathcal{B} : a(z)b(w) \text{ is regular for all } b \in \mathcal{B}\}$$

Since $\alpha(z)\alpha(w) \sim k/(z-w)^2$, the only fields with trivially regular OPE against all of~$\mathcal{B}$ are the scalars $\mathbb{C} \cdot \mathbf{1}$.  Hence:
$$HH^0(\mathcal{B}) = \mathbb{C}.$$

\subsubsection{\texorpdfstring{Degree 1: $HH^1(\mathcal{B})$}{Degree 1: HH 1(B)}}

$$HH^1 = \text{Der}(\mathcal{B})/\text{Inn}(\mathcal{B})$$
where $\text{Der}(\mathcal{B})$ is the space of derivations and $\text{Inn}(\mathcal{B})$ the inner derivations.  We show $\text{Der}(\mathcal{B}) = 0$ (hence a fortiori $HH^1 = 0$).

A derivation satisfies:
$$D(\alpha(z_1)\alpha(z_2)) = D(\alpha(z_1))\alpha(z_2) + \alpha(z_1)D(\alpha(z_2))$$

Using the OPE $\alpha \times \alpha \sim k/(z_1-z_2)^2$:
$$D\left(\frac{k}{(z_1-z_2)^2}\right) = \frac{D(\alpha)(z_1) \alpha(z_2) + \alpha(z_1)D(\alpha)(z_2)}{(z_1-z_2)^2}$$

The left side is a $c$-number (no field dependence), while the right side has field
dependence unless $D(\alpha) = 0$.

Therefore: $D = 0$.

$$HH^1(\mathcal{B}) = 0.$$

\subsubsection{\texorpdfstring{Degree 2: $HH^2(\mathcal{B})$}{Degree 2: HH 2(B)}}

$HH^2$ classifies deformations of the chiral algebra structure.

A deformation is given by modifying the OPE:
$$\alpha(z_1) \times_t \alpha(z_2) = \alpha(z_1)\alpha(z_2) + t \cdot \phi(z_1, z_2) + O(t^2)$$

where $\phi$ is a 2-cocycle.

\textbf{Cocycle condition:} Associativity to first order in $t$:
$$(\alpha \times_t \alpha) \times_t \alpha = \alpha \times_t (\alpha \times_t \alpha)$$

This imposes:
$$d\phi = 0$$
(cohomological condition)

\textbf{Trivial cocycles:} If $\phi = d\psi$ for some $\psi$, the deformation is trivial
(comes from a redefinition of $\alpha$).

\textbf{Nontrivial deformations:} The only independent deformation is changing the level $k$:
$$\alpha(z_1)\alpha(z_2) \sim \frac{k + t}{(z_1-z_2)^2}$$

This gives a 1-dimensional space:
$$HH^2(\mathcal{B}) = \mathbb{C} \cdot [\kappa].$$

where $[k]$ is the cohomology class of the level.

\subsubsection{\texorpdfstring{Higher Degrees: $HH^n(\mathcal{B})$ for $n \geq 3$}{Higher Degrees: HH n(B) for n >= 3}}

For the Heisenberg algebra (free boson), all higher Hochschild cohomology vanishes:
$$HH^n(\mathcal{B}) = 0 \quad \text{for } n \geq 3$$

\textbf{Proof.} The Heisenberg algebra $\mathcal{B}$ is Koszul: its bar construction $\bar{B}(\mathcal{B})$ is quasi-isomorphic to $\mathrm{CE}(\mathfrak{h}_{-k})$ (Theorem~\ref{thm:heisenberg-bar}), which is concentrated in bar degrees $\leq 2$ (the only generating relation is the quadratic OPE $\alpha(z)\alpha(w) \sim k/(z-w)^2$, involving exactly two fields).  The Koszul resolution therefore has the form:
$$\mathcal{B} \otimes \mathcal{B}^! \otimes \mathcal{B} \;\xrightarrow{d_2}\; \mathcal{B} \otimes (s\alpha) \otimes \mathcal{B} \;\xrightarrow{d_1}\; \mathcal{B} \otimes \mathcal{B}$$
with $\mathcal{B}^! = \mathrm{CE}(\mathfrak{h}_{-k})$ generated by a single element $s\alpha$ in bar degree~1 and its square in bar degree~2.  Applying $\mathrm{Hom}_{\mathcal{B}\text{-}\mathcal{B}}(-, \mathcal{B})$ to this length-2 resolution yields a cochain complex concentrated in degrees $0, 1, 2$.  Hence $HH^n(\mathcal{B}) = 0$ for $n \geq 3$.

\subsubsection{Summary for Heisenberg}

$$HH^*(\mathcal{B}) = \begin{cases}
\mathbb{C} & n = 0 \text{ (endomorphisms)}\\
0 & n = 1 \text{ (no derivations)}\\
\mathbb{C} & n = 2 \text{ (level deformation)}\\
0 & n \geq 3 \text{ (no higher structure)}
\end{cases}$$

As a graded algebra:
$$HH^*(\mathcal{B}) \cong \mathbb{C}[c]/(c^2) \cong \mathbb{C} \oplus \mathbb{C} \cdot c$$
where $c$ is a degree-2 class corresponding to the central charge (with $c^2 = 0$ since $HH^n = 0$ for $n \geq 3$).
\end{example}

\subsection{Explicit Computation: Free Fermion}

\begin{example}[Hochschild of Free Fermion - Complete]\label{ex:HH-fermion-complete}
For the free fermion $\mathcal{F}$ with fields $\psi(z), \psi^*(z)$ and OPE:
$$\psi(z_1)\psi^*(z_2) \sim \frac{1}{z_1 - z_2}$$

\subsubsection{\texorpdfstring{Degree 0: $HH^0(\mathcal{F})$}{Degree 0: HH 0(F)}}

Endomorphisms must preserve the fermionic OPE. By similar reasoning to the bosonic case:
$$HH^0(\mathcal{F}) = \mathbb{C}$$

\subsubsection{\texorpdfstring{Degree 1: $HH^1(\mathcal{F})$}{Degree 1: HH 1(F)}}

Derivations of the fermionic algebra. The Leibniz rule for fermions includes sign factors:
$$D(\psi \psi^*) = D(\psi)\psi^* + (-1)^{|\psi|} \psi D(\psi^*)$$

Since the OPE has no deformation parameters, we find:
$$HH^1(\mathcal{F}) = 0$$

\subsubsection{\texorpdfstring{Degree 2: $HH^2(\mathcal{F})$}{Degree 2: HH 2(F)}}

Deformations of the fermionic structure.  The free fermion $bc$ ghost system admits a one-parameter family of deformations parametrized by the spin $\lambda$:  the conformal weights $(h_b, h_c) = (\lambda, 1-\lambda)$ vary while the OPE $b(z)c(w) \sim 1/(z-w)$ is preserved.  This gives:
$$HH^2(\mathcal{F}) = \mathbb{C}$$

\subsubsection{Summary for Free Fermion}

$$HH^*(\mathcal{F}) = \begin{cases}
\mathbb{C} & n = 0\\
0 & n = 1\\
\mathbb{C} & n = 2\\
0 & n \geq 3
\end{cases}$$

As a graded algebra:
$$HH^*(\mathcal{F}) \cong \Lambda(c)$$
(exterior algebra on one generator of degree 2)
\end{example}

\subsection{Koszul Duality and HH* Pairing}

The Koszul duality for Hochschild cohomology (Theorem~\ref{thm:main-koszul-hoch}\label{thm:koszul-duality-HH}) is now verified on examples.

\begin{verification}[Boson-Fermion Duality; \ClaimStatusProvedHere]\label{ver:boson-fermion-HH}
The Koszul pair ($bc$ ghost system $\mathcal{F}_{bc}$, $\beta\gamma$ system $\mathcal{BG}$) satisfies (where the Hochschild cohomology of $\mathcal{F}_{bc}$ coincides with that of $\mathcal{B}$ in these degrees):
\begin{align*}
HH^0(\mathcal{B}) = \mathbb{C} &\xleftrightarrow{\text{dual}} HH^2(\mathcal{BG})^* = \mathbb{C}\\
HH^1(\mathcal{B}) = 0 &\xleftrightarrow{\text{dual}} HH^1(\mathcal{BG})^* = 0\\
HH^2(\mathcal{B}) = \mathbb{C} &\xleftrightarrow{\text{dual}} HH^0(\mathcal{BG})^* = \mathbb{C}
\end{align*}

This confirms the duality.
\end{verification}

\subsection{Comparison with Classical Hochschild Cohomology}

\begin{remark}[Chiral vs. Classical]\label{rem:chiral-vs-classical-hochschild}
Classical Hochschild cohomology (for associative algebras) is defined similarly:
$$HH^n_{\text{classical}}(A) = \text{Ext}^n_{A^e}(A, A)$$

\textbf{Key differences with chiral version:}

\begin{enumerate}
\item \textbf{Locality}: Chiral algebras have locality axiom (fields commute away from
coincident points), adding geometric structure

\item \textbf{Configuration spaces}: Chiral HH* involves integrals over configuration spaces,
while classical HH* is purely algebraic

\item \textbf{Differential forms}: Chiral cochains are tensored with differential forms,
encoding geometric information

\item \textbf{Dimensionality}: For vertex algebras, the cochains have an additional ``spatial''
direction from the curve $X$
\end{enumerate}

\textbf{Relation:} For the ``constant loops'' (fields independent of $z$), chiral HH* reduces
to classical HH*:
$$HH^n_{\text{chiral}}(\mathcal{A})|_{\text{constant}} \cong HH^n_{\text{classical}}(H^0(\mathcal{A}))$$
\end{remark}

\subsection{The Gerstenhaber Bracket from Configuration Spaces}

\begin{theorem}[Gerstenhaber Structure on HH*; \ClaimStatusProvedElsewhere]\label{thm:gerstenhaber-structure}
Hochschild cohomology carries a \textbf{Gerstenhaber bracket}:
$$[\cdot, \cdot]: HH^p(\mathcal{A}) \otimes HH^q(\mathcal{A}) \to HH^{p+q-1}(\mathcal{A})$$

making $HH^*(\mathcal{A})$ into a graded Lie algebra (with appropriate degree shift).
\end{theorem}

\begin{construction}[Geometric Realization of Bracket]\label{const:gerstenhaber-bracket}
The Gerstenhaber bracket has a beautiful geometric interpretation via configuration spaces.

For $f \in HH^p$ and $g \in HH^q$, represented as:
\begin{align*}
f &\in \Gamma(\overline{C}_{p+2}(X), \ldots)\\
g &\in \Gamma(\overline{C}_{q+2}(X), \ldots)
\end{align*}

The bracket $[f, g]$ is constructed by:
\begin{enumerate}
\item \textbf{Diagonal insertion}: Insert configuration of $f$ ``inside'' configuration of $g$
\item \textbf{Summation}: Sum over all possible insertion points
\item \textbf{Residue}: Extract the coefficient of singular terms
\end{enumerate}

Explicitly:
$$[f, g] = \sum_{i=1}^{q+1} (-1)^{\epsilon_i} \text{Res}_{z_0 \to z_i}\left[f(z_0, z_1, \ldots, z_p) 
\cdot g(\ldots, z_i, \ldots)\right]$$

where the residue extracts the collision behavior as one configuration approaches another.
\end{construction}

\begin{example}[Gerstenhaber Bracket for Heisenberg]\label{ex:gerstenhaber-heisenberg}
For $\mathcal{B}$ (Heisenberg), $HH^2 = \mathbb{C} \cdot [k]$ (level class).

The bracket:
$$[[k], [k]] = [k, k]$$

must have degree $2 + 2 - 1 = 3$. But $HH^3(\mathcal{B}) = 0$, so:
$$[[k], [k]] = 0$$

This reflects that the level is a \textbf{central element} in the Lie algebra structure
(it commutes with everything).
\end{example}

\subsection{\texorpdfstring{Higher Structure: $L_\infty$ Operations}{Higher Structure: L infty Operations}}

\begin{remark}[Beyond Gerstenhaber: Full $L_\infty$]\label{rem:L-infty-structure}
The configuration space geometry actually encodes a full \textbf{$L_\infty$ structure} on
$HH^*(\mathcal{A})$, not just the binary bracket.

\textbf{$L_\infty$ operations:}
$$\ell_n: HH^{*}(\mathcal{A})^{\otimes n} \to HH^{*+2-n}(\mathcal{A})$$

for all $n \geq 1$, satisfying higher Jacobi identities.

\textbf{Geometric realization:} Each $\ell_n$ comes from an integral over configuration
spaces with $n$ ``inputs'' and one ``output'':
$$\ell_n(f_1, \ldots, f_n) = \int_{\text{Config}_{n,1}} f_1 \wedge \cdots \wedge f_n \wedge K_n$$

where $K_n$ is a universal kernel (the ``associahedron measure'').

This $L_\infty$ structure governs the \textbf{deformation theory} of the chiral algebra completely.
\end{remark}

\subsection{Computational Strategy}

\begin{remark}[Computing $HH^*$ in practice]\label{rem:compute-HH}
To compute $HH^n_{\mathrm{chiral}}(\mathcal{A})$ for a chiral algebra given by generators and OPE data, one proceeds as follows: (1) construct the bar complex $\bar{B}^n(\mathcal{A})$ using configuration space integrals; (2) build the cobar complex $\Omega(\bar{B}(\mathcal{A}))$ by inserting distributional duals; (3) form the Hom complex $\operatorname{Hom}(\Omega(\bar{B}(\mathcal{A})), \mathcal{A})$; (4) compute cohomology $\ker(d)/\operatorname{im}(d)$ at each degree. For Koszul algebras, the quasi-isomorphism $\bar{B}(\mathcal{A}) \simeq \mathcal{A}^!$ reduces the computation to lower-dimensional complexes.
\end{remark}

\subsection{Summary and Outlook}

\begin{conclusion}
The bar-cobar approach to Hochschild cohomology provides:

\begin{enumerate}
\item \textbf{Explicit formulas}: All HH* groups computable via configuration space integrals

\item \textbf{Geometric understanding}: HH* arises from topology of configuration spaces

\item \textbf{Rich structure}: Gerstenhaber bracket and $L_\infty$ operations from geometry

\item \textbf{Koszul duality}: Perfect pairing $HH^n(\mathcal{A}) \leftrightarrow HH^{2-n}(\mathcal{A}^!)$

\item \textbf{Deformation theory}: $HH^2$ controls deformations, obstructions in $HH^3$
\end{enumerate}

This computational framework is applied throughout the manuscript to understand:
\begin{itemize}
\item Quantum corrections (obstruction theory)
\item W-algebras (screening charge cohomology)
\item Level-rank duality (Koszul pairing)
\item Higher genus contributions (moduli space cohomology)
\end{itemize}
\end{conclusion}
